\documentclass[11pt]{article}
% ===========================================================================
%  Shared preamble for all MPM2D worksheets — input right after \documentclass.
%  (Files beginning with "_" are skipped by mpm2d-worksheet-publish.mjs.)
% ===========================================================================
\usepackage[letterpaper,top=2.2cm,bottom=1.9cm,left=1.6cm,right=1.6cm,headheight=28pt,footskip=24pt]{geometry}
\usepackage{amsmath,amssymb}
\usepackage[T1]{fontenc}
\usepackage[utf8]{inputenc}
\usepackage{lmodern}
\usepackage{xcolor}
\usepackage[most]{tcolorbox}
\usepackage{enumitem}
\usepackage{multicol}
\usepackage{fancyhdr}
\usepackage{tikz}
\usetikzlibrary{calc}
\usepackage{pgfplots}
\pgfplotsset{compat=1.18}

\definecolor{exblue}{HTML}{4A90E2}
\definecolor{exbg}{HTML}{E6F3FF}
\definecolor{qorange}{HTML}{E69138}
\definecolor{qbg}{HTML}{FFF7CC}
\definecolor{keygray}{HTML}{475569}
\definecolor{keybg}{HTML}{F1F5F9}
\definecolor{iagreen}{HTML}{14653B}
\definecolor{titleblue}{HTML}{1D4ED8}
% Rotating palette so each worked example gets its own colour.
\definecolor{ex0}{HTML}{2563A0}\definecolor{ex1}{HTML}{3B7D3B}\definecolor{ex2}{HTML}{8A5A00}
\definecolor{ex3}{HTML}{A3327A}\definecolor{ex4}{HTML}{6D28A3}\definecolor{ex5}{HTML}{B08900}
\definecolor{ex6}{HTML}{0E7490}\definecolor{ex7}{HTML}{9A3412}\definecolor{ex8}{HTML}{15803D}

\pagestyle{fancy}
\fancyhf{}
\fancyhead[L]{\footnotesize NAME: \underline{\hspace{3.4cm}}}
\fancyhead[C]{\textbf{Integration Academy}\\[-2pt]{\scriptsize Math Simplified \textperiodcentered{} Success Amplified}}
\fancyhead[R]{\footnotesize MPM2D}
\fancyfoot[L]{\scriptsize dr.merie@IntegrationAcademy.ca}
\fancyfoot[C]{\scriptsize\thepage}
\fancyfoot[R]{\scriptsize IntegrationAcademy.ca}
\renewcommand{\headrulewidth}{1.2pt}
\renewcommand{\footrulewidth}{0.4pt}
\setlength{\parindent}{0pt}
\setlength{\parskip}{4pt}

% Examples are NOT breakable, so each example + its solution stays on one page.
\newtcolorbox{example}[2]{colframe=#1,colback=#1!7,coltitle=white,colbacktitle=#1,fonttitle=\bfseries,title={#2},arc=3pt,boxrule=1pt}
\newtcolorbox{qbox}[1]{colframe=qorange,colback=qbg,coltitle=white,colbacktitle=qorange,fonttitle=\bfseries,title={#1},arc=3pt,breakable,boxrule=1pt}
\newtcolorbox{keybox}{colframe=keygray,colback=keybg,coltitle=white,colbacktitle=keygray,fonttitle=\bfseries,title={$\checkmark$\ Answer Key},arc=3pt,breakable,boxrule=1pt}
\newtcolorbox{recapbox}{colframe=keygray,colback=keybg,coltitle=white,colbacktitle=keygray,fonttitle=\bfseries,title={Question Recap},arc=3pt,breakable,boxrule=1pt}
\newtcolorbox{ideabox}{colframe=titleblue,colback=titleblue!6,coltitle=white,colbacktitle=titleblue,fonttitle=\bfseries,title={The Big Ideas},arc=3pt,breakable,boxrule=1.2pt}
\newtcolorbox{learnbox}{colframe=iagreen,colback=iagreen!5,coltitle=white,colbacktitle=iagreen,fonttitle=\bfseries,title={Core Concepts},arc=3pt,breakable,boxrule=1.2pt}
\newcommand{\lh}[1]{\par\smallskip{\bfseries\color{iagreen}#1}\par\nobreak\smallskip}

\newcommand{\soln}{\par\smallskip\textit{\textbf{Solution.}}\ }
\newcommand{\workspace}[1]{\par\vspace{3pt}\fcolorbox{black!30}{white}{\begin{minipage}[t][#1][t]{\dimexpr\linewidth-2\fboxsep\relax}\ \end{minipage}}\par}
\newcommand{\sech}[1]{\par\vspace{8pt}{\Large\bfseries\color{iagreen}#1}\par\nobreak\vspace{2pt}{\color{black!20}\hrule height1pt}\vspace{6pt}}

% A compact coordinate plot sized for an example box: \eplot{xmin}{xmax}{ymin}{ymax}{body}
\newcommand{\eplot}[5]{\begin{center}\begin{tikzpicture}
\begin{axis}[width=6.8cm,height=4.4cm,axis lines=middle,xlabel={\tiny$x$},ylabel={\tiny$y$},
grid=both,grid style={gray!16},xmin=#1,xmax=#2,ymin=#3,ymax=#4,
xtick distance=2,ytick distance=2,tick label style={font=\tiny},axis line style={-},samples=120]
#5
\end{axis}\end{tikzpicture}\end{center}}

% Same as \eplot but with its own tick spacing: \eplotx{xmin}{xmax}{ymin}{ymax}{xtick}{ytick}{body}
\newcommand{\eplotx}[7]{\begin{center}\begin{tikzpicture}
\begin{axis}[width=8.4cm,height=5.4cm,axis lines=middle,xlabel={\tiny$x$},ylabel={\tiny$y$},
grid=both,grid style={gray!16},xmin=#1,xmax=#2,ymin=#3,ymax=#4,
xtick distance=#5,ytick distance=#6,tick label style={font=\tiny},axis line style={-},samples=120]
#7
\end{axis}\end{tikzpicture}\end{center}}

% A sine-curve plot (degrees on x, 0..360): \splot{ymin}{ymax}{body}
\newcommand{\splot}[3]{\begin{center}\begin{tikzpicture}
\begin{axis}[width=8.6cm,height=4.6cm,axis lines=middle,xlabel={\tiny$x\,(^\circ)$},ylabel={\tiny$y$},
grid=both,grid style={gray!16},xmin=0,xmax=360,ymin=#1,ymax=#2,
xtick distance=90,ytick distance=1,tick label style={font=\tiny},axis line style={-},samples=180]
#3
\end{axis}\end{tikzpicture}\end{center}}

% A small coordinate plot: \plot{xmin}{xmax}{ymin}{ymax}{body}
\newcommand{\plot}[5]{\begin{center}\begin{tikzpicture}
\begin{axis}[width=8.4cm,height=6cm,axis lines=middle,xlabel={\scriptsize$x$},ylabel={\scriptsize$y$},
grid=both,grid style={gray!16},xmin=#1,xmax=#2,ymin=#3,ymax=#4,
xtick distance=2,ytick distance=2,tick label style={font=\tiny},axis line style={-}]
#5
\end{axis}\end{tikzpicture}\end{center}}

% Right triangle (angle theta at bottom-left, right angle at bottom-right):
%   \rtri{drawBase}{drawHeight}{oppLabel}{adjLabel}{hypLabel}
\newcommand{\rtri}[5]{\begin{center}\begin{tikzpicture}[scale=0.85]
\coordinate (A) at (0,0);\coordinate (B) at (#1,0);\coordinate (C) at (#1,#2);
\draw[thick,exblue] (A)--(B)--(C)--cycle;
\draw[black] ($(B)+(-0.32,0)$)--($(B)+(-0.32,0.32)$)--($(B)+(0,0.32)$);
\node[below] at ($(A)!0.5!(B)$){#4};
\node[right] at ($(B)!0.5!(C)$){#3};
\node[above left] at ($(A)!0.5!(C)$){#5};
\node at (0.7,0.22){$\theta$};
\end{tikzpicture}\end{center}}

% General (oblique) triangle with vertex labels and opposite-side labels:
%   \gtri{Alabel}{Blabel}{Clabel}{aSide}{bSide}{cSide}
\newcommand{\gtri}[6]{\begin{center}\begin{tikzpicture}[scale=0.85]
\coordinate (A) at (0,0);\coordinate (B) at (5,0);\coordinate (C) at (1.6,3);
\draw[thick,exblue] (A)--(B)--(C)--cycle;
\node[below left] at (A){#1};\node[below right] at (B){#2};\node[above] at (C){#3};
\node[below] at ($(A)!0.5!(B)$){#6};
\node[right] at ($(B)!0.5!(C)$){#4};
\node[left] at ($(A)!0.5!(C)$){#5};
\end{tikzpicture}\end{center}}

\fancyhead[R]{\footnotesize MHF4U \textperiodcentered{} 4.2}
\begin{document}

\begin{center}
{\LARGE\bfseries\color{titleblue}Worksheet 4.2 --- Graphs of Logarithmic Functions}\\[2pt]
{\small Grade 12 Advanced Functions (MHF4U) \textperiodcentered{} Unit 4: Exponential and Logarithmic Functions}
\end{center}

\begin{ideabox}
$y=\log_b x$ is the reflection of $y=b^x$ in $y=x$: vertical asymptote $x=0$, domain $x>0$, range all reals, through $(1,0)$. In every step write the \textbf{whole function}, e.g.\ \emph{the argument of $y=\log(x-2)$ must be positive}.
\begin{itemize}[leftmargin=*,itemsep=2pt]
  \item VA at $x=0$; domain $x>0$ (the argument must be positive); range all reals.
  \item Passes through $(1,0)$ and $(b,1)$; increasing for $b>1$.
  \item For $y=a\log_b\big(k(x-d)\big)+c$ the VA is $x=d$ and points map by $(x,y)\to\left(\tfrac xk+d,\ ay+c\right)$.
\end{itemize}
\end{ideabox}

\begin{learnbox}
\lh{Inverse of the exponential}$y=\log_b x$ is the reflection of $y=b^x$ in the line $y=x$, with a \textbf{vertical} asymptote at $x=0$.
\lh{Domain and range}A logarithmic function is defined only for $x>0$, and its range is all real numbers — the mirror of the exponential.
\lh{Transformations}The framework $y=a\log_b\!\left(k(x-d)\right)+c$ shifts and stretches the curve and moves the vertical asymptote to $x=d$.
\lh{Writing solutions}Write the \textbf{whole function} in every step. To graph $y=a\log_b\!\big(k(x-d)\big)+c$: draw the asymptote $x=d$, map the parent's key points $\left(\tfrac1b,-1\right),(1,0),(b,1)$ with $(x,y)\to\left(\tfrac xk+d,\ ay+c\right)$, then sketch.
\end{learnbox}

\sech{Worked Examples}

\begin{example}{ex0}{Example 1 --- The log graph}
Sketch $y=\log_4 x$.\soln \textbf{Step 1 (vertical asymptote and domain):} The argument of $y=\log_4x$ must be positive, so the domain is $x>0$ and the VA is $x=0$.

\textbf{Step 2 (key points):} From $4^y=x$: $y=-1\Rightarrow x=\tfrac14$; $y=0\Rightarrow x=1$; $y=1\Rightarrow x=4$; $y=2\Rightarrow x=16$. So $\left(\tfrac14,-1\right),(1,0),(4,1),(16,2)$ are on $y=\log_4x$.

\textbf{Step 3 (shape):} The base $4>1$, so the curve is increasing, with range all reals.

\textbf{Answer:} $y=\log_4x$: VA $x=0$, increasing, through $(1,0)$ and $(4,1)$.

\eplotx{0}{18}{-3}{3}{4}{1}{\addplot[exblue,very thick,domain=0.05:18,samples=120]{ln(x)/ln(4)};\draw[dashed,gray] (axis cs:0,-3)--(axis cs:0,3);\addplot[red,only marks,mark=*,mark size=1.6pt] coordinates {(0.25,-1) (1,0) (4,1) (16,2)};\node[font=\tiny,right] at (axis cs:0.25,-1) {$(1/4,-1)$};\node[font=\tiny,below right] at (axis cs:1,0) {$(1,0)$};\node[font=\tiny,below right] at (axis cs:4,1) {$(4,1)$};\node[font=\tiny,below left] at (axis cs:16,2) {$(16,2)$};}
\end{example}

\begin{example}{ex1}{Example 2 --- Domain}
Domain of $y=\log x$?\soln \textbf{Step 1 (positive argument):} The argument of $y=\log x$ must be positive: $x>0$.

\textbf{Step 2 (test values):} $y=\log5$ is defined, but $y=\log0$ and $y=\log(-3)$ are not.

\textbf{Answer:} the domain of $y=\log x$ is $x>0$.

\eplotx{0}{12}{-2}{3}{2}{1}{\addplot[exblue,very thick,domain=0.05:12,samples=120]{log10(x)};\draw[dashed,gray] (axis cs:0,-2)--(axis cs:0,3);\addplot[red,only marks,mark=*,mark size=1.6pt] coordinates {(1,0) (10,1)};\node[font=\tiny,below right] at (axis cs:1,0) {$(1,0)$};\node[font=\tiny,below right] at (axis cs:10,1) {$(10,1)$};}
\end{example}

\begin{example}{ex2}{Example 3 --- Inverse pair}
How are $y=5^x$ and $y=\log_5 x$ related?\soln \textbf{Step 1 (swap x and y):} Start from $y=5^x$ and swap $x$ and $y$: $x=5^y\iff y=\log_5x$.

\textbf{Step 2 (points swap):} A point $(a,b)$ on $y=5^x$ becomes $(b,a)$ on $y=\log_5x$: $(0,1)\to(1,0)$ and $(1,5)\to(5,1)$.

\textbf{Step 3 (reflect):} The graphs are mirror images across the line $y=x$; the horizontal asymptote $y=0$ of $y=5^x$ becomes the vertical asymptote $x=0$ of $y=\log_5x$.

\textbf{Answer:} $y=5^x$ and $y=\log_5x$ are inverses (reflections in $y=x$).

\eplotx{-2}{6}{-2}{6}{2}{2}{\addplot[exblue,very thick,domain=-2:1.1,samples=80]{5^x};\addplot[exblue,very thick,domain=0.05:6,samples=120]{ln(x)/ln(5)};\addplot[gray,dashed,very thick,domain=-2:6,samples=120]{x};\draw[dashed,gray] (axis cs:0,-2)--(axis cs:0,6);\addplot[red,only marks,mark=*,mark size=1.6pt] coordinates {(0,1) (1,5) (1,0) (5,1)};\node[font=\tiny,above left] at (axis cs:0,1) {$(0,1)$};\node[font=\tiny,right] at (axis cs:1,5) {$(1,5)$};\node[font=\tiny,below right] at (axis cs:1,0) {$(1,0)$};\node[font=\tiny,below right] at (axis cs:5,1) {$(5,1)$};}
\end{example}

\begin{example}{ex3}{Example 4 --- Vertical asymptote}
VA of $y=\log x$?\soln \textbf{Step 1 (look near zero):} For $y=\log x$: $\log0.1=-1$, $\log0.001=-3$, $\log0.000001=-6$. As $x\to0^+$, $y\to-\infty$.

\textbf{Step 2 (asymptote):} The curve $y=\log x$ falls without bound but never reaches the y-axis.

\textbf{Answer:} the vertical asymptote of $y=\log x$ is $x=0$.
\end{example}

\begin{example}{ex4}{Example 5 --- Key point}
What point is on every graph $y=\log_b x$?\soln \textbf{Step 1 (log of 1):} For $y=\log_bx$ set $x=1$: $\log_b1=y\iff b^y=1$, so $y=0$ for every base $b$.

\textbf{Step 2 (also (b, 1)):} Since $\log_bb=1$, the point $(b,1)$ is also on $y=\log_bx$: $(2,1)$, $(5,1)$, $(10,1)$.

\textbf{Answer:} $(1,0)$, whatever the base.

\eplotx{0}{12}{-2}{3}{2}{1}{\addplot[exblue,very thick,domain=0.05:12,samples=120]{ln(x)/ln(2)};\addplot[qorange,very thick,domain=0.05:12,samples=120]{ln(x)/ln(5)};\addplot[red,very thick,domain=0.05:12,samples=120]{log10(x)};\draw[dashed,gray] (axis cs:0,-2)--(axis cs:0,3);\addplot[red,only marks,mark=*,mark size=1.6pt] coordinates {(1,0)};\node[font=\tiny,below right] at (axis cs:1,0) {$(1,0)$};}
\end{example}

\begin{example}{ex5}{Example 6 --- Transformation}
VA of $y=\log(x-2)$?\soln \textbf{Step 1 (positive argument):} The argument of $y=\log(x-2)$ must be positive: $x-2>0\Rightarrow x>2$.

\textbf{Step 2 (asymptote):} As $x\to2^+$, $y=\log(x-2)\to-\infty$, so the VA is $x=2$: the graph of $y=\log x$ has moved right $2$.

\textbf{Step 3 (points):} $(1,0)\to(3,0)$ and $(10,1)\to(12,1)$. Check: $\log(3-2)=\log1=0$ $\checkmark$.

\textbf{Answer:} the vertical asymptote of $y=\log(x-2)$ is $x=2$.

\eplotx{0}{14}{-3}{3}{2}{1}{\addplot[exblue,very thick,domain=2.05:14,samples=120]{log10(x-2)};\draw[dashed,gray] (axis cs:2,-3)--(axis cs:2,3);\addplot[red,only marks,mark=*,mark size=1.6pt] coordinates {(3,0) (12,1)};\node[font=\tiny,below right] at (axis cs:3,0) {$(3,0)$};\node[font=\tiny,below left] at (axis cs:12,1) {$(12,1)$};}
\end{example}

\begin{example}{ex6}{Example 7 --- Range}
Range of $y=\log x$?\soln \textbf{Step 1 (can y be any number?):} For $y=\log x$, every real $y$ has an $x$ with $\log x=y$, namely $x=10^y$.

\textbf{Step 2 (examples):} $y=5$ comes from $x=10^5$; $y=-5$ comes from $x=10^{-5}$.

\textbf{Answer:} the range of $y=\log x$ is all real numbers.
\end{example}

\begin{example}{ex7}{Example 8 --- Inverse}
Inverse of $y=7^x$?\soln \textbf{Step 1 (swap x and y):} Start from $y=7^x$ and swap: $x=7^y$.

\textbf{Step 2 (logarithmic form):} $x=7^y$ is the same as $y=\log_7x$.

\textbf{Step 3 (check with a point):} $(1,7)$ on $y=7^x$ becomes $(7,1)$ on $y=\log_7x$, and $\log_77=1$ $\checkmark$.

\textbf{Answer:} the inverse of $y=7^x$ is $y=\log_7x$.
\end{example}

\begin{example}{ex8}{Example 9 --- Transformation}
VA of $y=\log(x+6)$?\soln \textbf{Step 1 (positive argument):} The argument of $y=\log(x+6)$ must be positive: $x+6>0\Rightarrow x>-6$.

\textbf{Step 2 (asymptote):} As $x\to-6^+$, $y=\log(x+6)\to-\infty$, so the VA is $x=-6$: the graph of $y=\log x$ has moved left $6$.

\textbf{Step 3 (points):} $(1,0)\to(-5,0)$ and $(10,1)\to(4,1)$. Check: $\log(4+6)=\log10=1$ $\checkmark$.

\textbf{Answer:} the vertical asymptote of $y=\log(x+6)$ is $x=-6$.

\eplotx{-8}{8}{-3}{3}{2}{1}{\addplot[exblue,very thick,domain=-5.95:8,samples=120]{log10(x+6)};\draw[dashed,gray] (axis cs:-6,-3)--(axis cs:-6,3);\addplot[red,only marks,mark=*,mark size=1.6pt] coordinates {(-5,0) (4,1)};\node[font=\tiny,below right] at (axis cs:-5,0) {$(-5,0)$};\node[font=\tiny,below right] at (axis cs:4,1) {$(4,1)$};}
\end{example}

\clearpage
\sech{Your Turn --- Show Your Work}
For each question, show all steps clearly in the space provided.

\begin{qbox}{Question 1}
Domain of $y=\log_6 x$?\workspace{2.4cm}
\end{qbox}

\begin{qbox}{Question 2}
VA of $y=\log(x+8)$?\workspace{2.4cm}
\end{qbox}

\begin{qbox}{Question 3}
What is $\log_b b$?\workspace{2.4cm}
\end{qbox}

\begin{qbox}{Question 4}
Is $y=\log_7 x$ increasing or decreasing?\workspace{2.4cm}
\end{qbox}

\begin{qbox}{Question 5}
Inverse of $y=4^x$?\workspace{2.4cm}
\end{qbox}

\begin{qbox}{Question 6}
Range of $y=\log x$?\workspace{2.4cm}
\end{qbox}

\begin{qbox}{Question 7}
VA of $y=\log(x-5)$?\workspace{2.4cm}
\end{qbox}

\begin{qbox}{Question 8}
Domain of $y=\log(x-1)$?\workspace{2.4cm}
\end{qbox}

\begin{qbox}{Question 9}
Through which point does $y=\log_9 x$ pass on the x-axis?\workspace{2.4cm}
\end{qbox}

\begin{qbox}{Question 10}
Inverse of $y=10^x$?\workspace{2.4cm}
\end{qbox}

\begin{qbox}{Question 11}
VA of $y=\log(x)+4$?\workspace{2.4cm}
\end{qbox}

\begin{qbox}{Question 12}
As $x\to0^+$, what does $y=\log x$ do?\workspace{2.4cm}
\end{qbox}

\begin{qbox}{Question 13 --- Challenge}
State the domain, range, VA and a key point of $y=\log_2(x-1)$.\workspace{3.5cm}
\end{qbox}

\clearpage
\sech{Question Recap \& Answer Key}

\begin{recapbox}
\begin{multicols}{2}
\small
\begin{enumerate}[leftmargin=*,itemsep=3pt]
  \item Domain of $y=\log_6 x$?
  \item VA of $y=\log(x+8)$?
  \item What is $\log_b b$?
  \item Is $y=\log_7 x$ increasing or decreasing?
  \item Inverse of $y=4^x$?
  \item Range of $y=\log x$?
  \item VA of $y=\log(x-5)$?
  \item Domain of $y=\log(x-1)$?
  \item Through which point does $y=\log_9 x$ pass on the x-axis?
  \item Inverse of $y=10^x$?
  \item VA of $y=\log(x)+4$?
  \item As $x\to0^+$, what does $y=\log x$ do?
  \item State the domain, range, VA and a key point of $y=\log_2(x-1)$.
\end{enumerate}
\end{multicols}
\end{recapbox}

\begin{keybox}
\begin{multicols}{2}
\small
\begin{enumerate}[leftmargin=*,itemsep=3pt]
  \item For $y=\log_6x$ the argument must be positive, so the domain is $x>0$
  \item For $y=\log(x+8)$: $x+8>0\Rightarrow x>-8$, so the VA is $x=-8$
  \item $\log_bb=y\iff b^y=b^1$, so $\log_bb=1$
  \item For $y=\log_7x$ the base $7>1$, so the function is increasing
  \item Swap $x$ and $y$ in $y=4^x$: $x=4^y\iff y=\log_4x$
  \item The range of $y=\log x$ is all reals (e.g.\ $\log10^5=5$ and $\log10^{-5}=-5$)
  \item For $y=\log(x-5)$: $x-5>0\Rightarrow x>5$, so the VA is $x=5$
  \item For $y=\log(x-1)$: $x-1>0\Rightarrow x>1$, so the domain is $x>1$
  \item Set $y=0$ in $y=\log_9x$: $\log_9x=0\Rightarrow x=9^0=1$, so the point is $(1,0)$
  \item Swap $x$ and $y$ in $y=10^x$: $x=10^y\iff y=\log x$ ($=\log_{10}x$)
  \item For $y=\log x+4$ the argument is still $x$, so $x>0$ and the VA is $x=0$ (the $+4$ only moves the graph up)
  \item For $y=\log x$: $\log0.01=-2$ and $\log0.0001=-4$, so $y\to-\infty$
  \item For $y=\log_2(x-1)$: $x-1>0$ gives domain $x>1$ and VA $x=1$; range all reals; key points $(1,0)\to(2,0)$ and $(2,1)\to(3,1)$, so $(2,0)$ and $(3,1)$ are on the graph
\end{enumerate}
\end{multicols}
\end{keybox}

\end{document}
